\scanpage{7}
Reste à vérifier que $\varphi$ est \uline{continue}, \deleted{i.e. que l'image inverse par $\varphi$ d'un ouvert est un ouvert. C'est vrai sur $U=Y-y$ par hyp. de récurrence [car alors $X|U$ a une solution de modules pour} et on sait que cel. provient du fait que $X\to Y$ est \deleted{\unsure{univers\textsuperscript{t}}} ouvert.

\deleted{\uline{On obtient} ainsi}

Si on ne suppose plus $f^{-1}(y)$ irréductible, \inserted{mais ses composantes irréd. ult irréd.} on \unsure{trouve|trouver} par la méthode précédente \deleted{\unsure{on voit}} \inserted{qu'il existe $r$} fonctions \unsure{monisantes|monisentes} $\varphi_i$ $(1\leq i\leq r)$ \deleted{pour} de $X/Y$, \deleted{\unsure{correspondant} \unsure{qui sont les seules fonctions}} une pour chaque comp. de $f^{-1}(y)$. Soit \deleted{\unsure{$=$}} $Y_i=Y$ $(i=1,\unsure{--|\ldots,n})$, et $Y'_i={}$\deleted{\unsure{$X$}}$(Y-y)$, \deleted{\unsure{les}} Soit \deleted{\unsure{$X$}} enfin $M'$ la solution des modules pour $X'=X|Y'$,
% slanted
\margin{``hyp. de récurrence''}
alors les \deleted{\unsure{$X_i$}} $\varphi_i:Y'\to X$ correspondent à des sections $\varphi'_i$ de $M'/Y'$, donc à des ouverts $M'_i$ de $M'$ isomorphes aux $Y'_i$.

Recollons \deleted{les} $M'$ et les $Y'_i$ \unsure{par} \deleted{\unsure{le}} $Y'_i\to M'_i$. On trouve un préschéma $M$ \deleted{sur} quasi fini et étale sur $Y$. L'application structurale $\psi':M'\to X'\times_{Y'}M'=X\unsure{\times}{}_YM'$ et les $\varphi_i:\unsure{Y_i|X_i}\to X\inserted{=Y_i\times_Y X}$ définissent \deleted{\unsure{$\varphi$}} $\psi:M\to X\times_Y M$ qui sera une application monisante (caractère local d'une \unsure{apl.} monisante) et \unsure{sera} une solution du Pb des modules pour $X,Y$.
% archive sheet number: 6
